% =============================================================================
\subsection{平均场约化、矩闭包与序参量}
\label{sec:theory-mf}
格点耦合使精确求解不可行，但\emph{平均场}约化可以让集体物理变得透明，并
揭示一个在单体图像中看不见的有序--无序平衡。记
\begin{equation}
  \bar J:=\sum_j J_{ij},\qquad
  m(t):=\E[x_i(t)],\qquad
  V(t):=\E[x_i^2(t)]-m^2(t),
  \label{eq:mf-defs}
\end{equation}
分别为总耦合权重、\emph{序参量}（平均意见，度量共识程度）与涨落（度量为
分歧程度）。把耦合项近似为
$\sum_j J_{ij}(x_j-x_i)\approx \bar J\,(m-x_i)$，并用高斯闭包
$\E[x^3]=m^3+3mV$、$\E[x^4]=m^4+6m^2V+3V^2$ 闭合矩层级，得到
\begin{align}
  \dot m &= (a-3bV)\,m-b\,m^{3}+A\cos(\omega t)+C,
  \label{eq:mf-m}\\
  \dot V &= 2\Bigl[a(m^{2}+V)-b\bigl(m^{4}+6m^{2}V+3V^{2}\bigr)-\bar J V\Bigr]
           \nonumber\\
         &\quad +2\Bigl[A\cos(\omega t)+C\Bigr]m+2D .
  \label{eq:mf-V}
\end{align}
注意耦合以\emph{差分}形式进入，故它在 $\dot m$ 中完全抵消，只在 $\dot V$
中留下 $-\bar JV$ 项。这一结构性事实支配着仿真的解释方式，值得逐条说明。

\paragraph{序参量与临界涨落。}
取 $A=C=0$ 并寻找有序定态（$m\neq0$，$\dot m=0$），\eqref{eq:mf-m} 给出
\begin{equation}
  m^{2}=\frac{a}{b}-3V .
  \label{eq:mf-order}
\end{equation}
即序参量被涨落\emph{抑制}：随着噪声抬升 $V$，共识程度 $|m|$ 下降，并在临界
涨落
\begin{equation}
  V_c=\frac{a}{3b}\qquad (a=b=1\Rightarrow V_c=1/3)
  \label{eq:mf-critical}
\end{equation}
处消失。当 $V>V_c$ 时有序解不再存在，对称的无序态 $m=0$ 是唯一的定态。这就是
模型的平均场有序--无序相变，是社会从众性的 Curie--Weiss 型相变在社会系统中的
对应物。社会学地读：\emph{当分歧（涨落）超过一个临界值时，共识在数学上不再
可能存在}——不是难以达成，而是不存在。

\paragraph{耦合的作用。}
\eqref{eq:mf-V} 表明耦合以 $-\bar J V$ 进入方差方程：社会影响\emph{抑制}
分歧，因而通过 \eqref{eq:mf-order} \emph{促进}共识。等价地，强耦合使有序分支
对噪声稳定。这正是图~\ref{fig:snap} 快照中所见相干畴背后的机制：耦合把 $V$
压得足够低，使 $V<V_c$ 得以维持，从而支撑起一个大的 $|m|$。这也给出一个可检验
的预言：增大耦合强度 $\bar J$ 会把相变——以及随之而来的共振峰——推向更大的
$D$。

\paragraph{与共振的联系。}
平均场图像澄清了共振为什么必须位于\emph{中等}噪声处。在 $D$ 小时 $V$ 很小，
$m$ 被锁定在 $m\approx\pm\sqrt{a/b}$：群体是有序的但被\emph{冻结}，无法跟随
微弱的驱动，因此响应中相干（被驱动）的部分可以忽略。在 $D$ 大时 $V$ 超过
$V_c$，有序态被破坏，$m\to0$：群体是无序的，同样不携带相干信号。只有在
中等 $V$ 处——既非冻结亦非无序——群体才既有足够的秩序以集体响应，又有足够的
流动性以发生切换，而这恰恰就是共振条件。下一小节的 Kramers 分析将把这一平衡
定量化。

\paragraph{关于闭包精度的说明。}
高斯闭包仅在 $\alpha=2$ 且处于线性（单阱）极限时才是精确的；对重尾噪声，
真实的增量分布具有幂律尾，闭包会低估大偏离的概率。因此我们仅把
\eqref{eq:mf-m}--\eqref{eq:mf-critical} 用于定性解释，定量结论依赖直接仿真。

